\section{E}

% empfehlen (to recommend)
\begin{verbentry}{
  style = {default},
  toc = {empfehlen (to recommend)},
  name = {empfehlen},
  english = {to recommend},
  type = {irregular},
  auxiliary = {haben}
}


\vspace{5pt}
\hrule
\vspace{5pt}

\begin{tabular}{rl}
\multicolumn{2}{l}{\textbf{PRÄSENS}}\\
ich & \textbf{empfehle}\\
du & \textbf{empfiehlst}\\
er/sie/es & \textbf{empfiehlt}\\
wir & \textbf{empfehlen}\\
ihr & \textbf{empfehlt}\\
sie/Sie & \textbf{empfehlen}\\
\end{tabular}
\hfill
\begin{tabular}{rl}
\multicolumn{2}{l}{\textbf{PRÄTERITUM}}\\
ich & \textbf{empfahl}\\
du & \textbf{empfahlst}\\
er/sie/es & \textbf{empfahl}\\
wir & \textbf{empfahlen}\\
ihr & \textbf{empfahlt}\\
sie/Sie & \textbf{empfahlen}\\
\end{tabular}
\hfill
\begin{tabular}{rl}
\multicolumn{2}{l}{\textbf{IMPERATIV}}\\
& \textbf{empfiehl (du)!}\\
& \textbf{empfehlen wir!}\\
& \textbf{empfehlt (ihr)!}\\
& \textbf{empfehlen Sie!}\\
\end{tabular}

\vspace{10pt}

\begin{tabular}{rl}
\multicolumn{2}{l}{\textbf{PERFEKT}}\\
& haben + \textbf{empfohlen}\\
\end{tabular}
\hfill
\begin{tabular}{rl}
\multicolumn{2}{l}{\textbf{FUTURE I}}\\
& werden + \textbf{empfehlen}\\
\end{tabular}
\hfill
\begin{tabular}{rl}
\multicolumn{2}{l}{\textbf{PLUSQUAMPERFEKT}}\\
& hatten + \textbf{empfohlen}\\
\end{tabular}

\vspace{10pt}

\begin{tabular}{lll}
\multicolumn{3}{l}{\textbf{MIT PRÄPOSITIONEN}}\\
& \textbf{empfehlen für} + Akk & (to recommend for)\\
& \textbf{empfehlen an} + Akk & (to recommend to)\\
\end{tabular}
\hfill
\begin{tabular}{lll}
\multicolumn{3}{l}{\textbf{TRENNBARE VERBEN}}\\
& \textbf{---} & \\
\end{tabular}
\vspace{-5pt}
\end{verbentry}

% engagieren (to hire / engage)
\begin{verbentry}{
  style = {acc},
  toc = {engagieren (to hire / engage)},
  name = {engagieren},
  english = {to hire / engage},
  type = {regular, + Akkusativ},
  auxiliary = {haben}
}


\vspace{5pt}
\hrule
\vspace{5pt}

\begin{tabular}{rl}
\multicolumn{2}{l}{\textbf{PRÄSENS}}\\
ich & \textbf{engagiere}\\
du & \textbf{engagierst}\\
er/sie/es & \textbf{engagiert}\\
wir & \textbf{engagieren}\\
ihr & \textbf{engagiert}\\
sie/Sie & \textbf{engagieren}\\
\end{tabular}
\hfill
\begin{tabular}{rl}
\multicolumn{2}{l}{\textbf{PRÄTERITUM}}\\
ich & \textbf{engagierte}\\
du & \textbf{engagiertest}\\
er/sie/es & \textbf{engagierte}\\
wir & \textbf{engagierten}\\
ihr & \textbf{engagiertet}\\
sie/Sie & \textbf{engagierten}\\
\end{tabular}
\hfill
\begin{tabular}{rl}
\multicolumn{2}{l}{\textbf{IMPERATIV}}\\
& \textbf{engagiere (du)!}\\
& \textbf{engagieren wir!}\\
& \textbf{engagiert (ihr)!}\\
& \textbf{engagieren Sie!}\\
\end{tabular}

\vspace{10pt}

\begin{tabular}{rl}
\multicolumn{2}{l}{\textbf{PERFEKT}}\\
& haben + \textbf{engagiert}\\
\end{tabular}
\hfill
\begin{tabular}{rl}
\multicolumn{2}{l}{\textbf{FUTURE I}}\\
& werden + \textbf{engagieren}\\
\end{tabular}
\hfill
\begin{tabular}{rl}
\multicolumn{2}{l}{\textbf{PLUSQUAMPERFEKT}}\\
& hatten + \textbf{engagiert}\\
\end{tabular}

\vspace{10pt}

\begin{tabular}{lll}
\multicolumn{3}{l}{\textbf{MIT PRÄPOSITIONEN}}\\
& \textbf{---} & \\
\end{tabular}
\hfill
\begin{tabular}{lll}
\multicolumn{3}{l}{\textbf{TRENNBARE VERBEN}}\\
& \textbf{---} & \\
\end{tabular}

\vspace{-5pt}
\end{verbentry}

% entdecken (to discover)
\begin{verbentry}{
  style = {acc},
  toc = {entdecken (to discover)},
  name = {entdecken},
  english = {to discover},
  type = {regular, + Akkusativ},
  auxiliary = {haben}
}


\vspace{5pt}
\hrule
\vspace{5pt}

\begin{tabular}{rl}
\multicolumn{2}{l}{\textbf{PRÄSENS}}\\
ich & \textbf{entdecke}\\
du & \textbf{entdeckst}\\
er/sie/es & \textbf{entdeckt}\\
wir & \textbf{entdecken}\\
ihr & \textbf{entdeckt}\\
sie/Sie & \textbf{entdecken}\\
\end{tabular}
\hfill
\begin{tabular}{rl}
\multicolumn{2}{l}{\textbf{PRÄTERITUM}}\\
ich & \textbf{entdeckte}\\
du & \textbf{entdecktest}\\
er/sie/es & \textbf{entdeckte}\\
wir & \textbf{entdeckten}\\
ihr & \textbf{entdecktet}\\
sie/Sie & \textbf{entdeckten}\\
\end{tabular}
\hfill
\begin{tabular}{rl}
\multicolumn{2}{l}{\textbf{IMPERATIV}}\\
& \textbf{entdecke (du)!}\\
& \textbf{entdecken wir!}\\
& \textbf{entdeckt (ihr)!}\\
& \textbf{entdecken Sie!}\\
\end{tabular}

\vspace{10pt}

\begin{tabular}{rl}
\multicolumn{2}{l}{\textbf{PERFEKT}}\\
& haben + \textbf{entdeckt}\\
\end{tabular}
\hfill
\begin{tabular}{rl}
\multicolumn{2}{l}{\textbf{FUTURE I}}\\
& werden + \textbf{entdecken}\\
\end{tabular}
\hfill
\begin{tabular}{rl}
\multicolumn{2}{l}{\textbf{PLUSQUAMPERFEKT}}\\
& hatten + \textbf{entdeckt}\\
\end{tabular}

\vspace{10pt}

\begin{tabular}{lll}
\multicolumn{3}{l}{\textbf{MIT PRÄPOSITIONEN}}\\
& \textbf{---} & \\
\end{tabular}
\hfill
\begin{tabular}{lll}
\multicolumn{3}{l}{\textbf{TRENNBARE VERBEN}}\\
& \textbf{---} & \\
\end{tabular}

\vspace{-5pt}
\end{verbentry}

% entfernen (to remove)
\begin{verbentry}{
  style = {acc},
  toc = {entfernen (to remove)},
  name = {entfernen},
  english = {to remove},
  type = {regular, + Akkusativ},
  auxiliary = {haben}
}


\vspace{5pt}
\hrule
\vspace{5pt}

\begin{tabular}{rl}
\multicolumn{2}{l}{\textbf{PRÄSENS}}\\
ich & \textbf{entferne}\\
du & \textbf{entfernst}\\
er/sie/es & \textbf{entfernt}\\
wir & \textbf{entfernen}\\
ihr & \textbf{entfernt}\\
sie/Sie & \textbf{entfernen}\\
\end{tabular}
\hfill
\begin{tabular}{rl}
\multicolumn{2}{l}{\textbf{PRÄTERITUM}}\\
ich & \textbf{entfernte}\\
du & \textbf{entferntest}\\
er/sie/es & \textbf{entfernte}\\
wir & \textbf{entfernten}\\
ihr & \textbf{entferntet}\\
sie/Sie & \textbf{entfernten}\\
\end{tabular}
\hfill
\begin{tabular}{rl}
\multicolumn{2}{l}{\textbf{IMPERATIV}}\\
& \textbf{entferne (du)!}\\
& \textbf{entfernen wir!}\\
& \textbf{entfernt (ihr)!}\\
& \textbf{entfernen Sie!}\\
\end{tabular}

\vspace{10pt}

\begin{tabular}{rl}
\multicolumn{2}{l}{\textbf{PERFEKT}}\\
& haben + \textbf{entfernt}\\
\end{tabular}
\hfill
\begin{tabular}{rl}
\multicolumn{2}{l}{\textbf{FUTURE I}}\\
& werden + \textbf{entfernen}\\
\end{tabular}
\hfill
\begin{tabular}{rl}
\multicolumn{2}{l}{\textbf{PLUSQUAMPERFEKT}}\\
& hatten + \textbf{entfernt}\\
\end{tabular}

\vspace{10pt}

\begin{tabular}{lll}
\multicolumn{3}{l}{\textbf{MIT PRÄPOSITIONEN}}\\
& \textbf{entfernen aus} + Dat & (to remove from)\\
\end{tabular}
\hfill
\begin{tabular}{lll}
\multicolumn{3}{l}{\textbf{TRENNBARE VERBEN}}\\
& \textbf{---} & \\
\end{tabular}

\vspace{-5pt}
\end{verbentry}

% entlassen (to dismiss / to release)
\begin{verbentry}{
  style = {default},
  toc = {entlassen (to dismiss / to release)},
  name = {entlassen},
  english = {to dismiss / to release},
  type = {irregular},
  auxiliary = {haben}
}
 
    
     \vspace{5pt}

    \hrule
    
    \vspace{5pt}

  \begin{tabular}{rl}
     \multicolumn{2}{l}{\textbf{PRÄSENS} }\\
    ich & \textbf{entlasse}\\
    du & \textbf{entlässt}\\
    er/sie/es & \textbf{entlässt}\\
    wir & \textbf{entlassen}\\
    ihr & \textbf{entlasst}\\
    sie/Sie & \textbf{entlassen}\\
  \end{tabular}
  \hfill
     \begin{tabular}{rl}
    \multicolumn{2}{l}{\textbf{PRÄTERITUM} }\\
    ich & \textbf{entließ}\\
    du & \textbf{entließt}\\
    er/sie/es & \textbf{entließ}\\
    wir & \textbf{entließen}\\
    ihr & \textbf{entließt}\\
    sie/Sie & \textbf{entließen}\\
  \end{tabular}
  \hfill
  \begin{tabular}{rl}
     \multicolumn{2}{l}{\textbf{IMPERATIV} }\\
   & \textbf{entlass(e) (du)!}\\
   & \textbf{entlassen wir!}\\
   & \textbf{entlasst (ihr)!}\\
   & \textbf{entlassen Sie!}\\
   & \vspace{10pt}
  \end{tabular}

    \vspace{10pt}

 \begin{tabular}{rl}
     \multicolumn{2}{l}{\textbf{PERFEKT}}\\
   & haben + \textbf{entlassen}\\
  \end{tabular}
\hfill
   \begin{tabular}{rl}
    \multicolumn{2}{l}{\textbf{FUTURE I} }\\
   & werden + \textbf{entlassen}\\
  \end{tabular}
\hfill
   \begin{tabular}{rl}
      \multicolumn{2}{l}{\textbf{PLUSQUAMPERFEKT}}\\
   & hatten + \textbf{entlassen}\\
  \end{tabular}

   \vspace{10pt}

   \begin{tabular}{lll}
      \multicolumn{3}{l}{\textbf{MIT PRÄPOSITIONEN}}\\
& \textbf{entlassen aus} + Dat & (to release from / dismiss from)\\
& \textbf{entlassen wegen} + Gen & (to dismiss because of)\\
\end{tabular}
  \hfill
   \begin{tabular}{lll}
\multicolumn{3}{l}{\textbf{TRENNBARE VERBEN}}\\
& \textbf{---} & \\
\end{tabular}

  \vspace{-5pt}
\end{verbentry}

% entscheiden (to decide)
\begin{verbentry}{
  style = {acc},
  toc = {entscheiden (to decide)},
  name = {entscheiden},
  english = {to decide},
  type = {irregular, + Akkusativ},
  auxiliary = {haben}
}
 
    
     \vspace{5pt}

    \hrule
    
    \vspace{5pt}

  \begin{tabular}{rl}
     \multicolumn{2}{l}{\textbf{PRÄSENS} }\\
    ich & \textbf{entscheide}\\
    du & \textbf{entscheidest}\\
    er/sie/es & \textbf{entscheidet}\\
    wir & \textbf{entscheiden}\\
    ihr & \textbf{entscheidet}\\
    sie/Sie & \textbf{entscheiden}\\
  \end{tabular}
  \hfill
     \begin{tabular}{rl}
    \multicolumn{2}{l}{\textbf{PRÄTERITUM} }\\
    ich & \textbf{entschied}\\
    du & \textbf{entschiedst}\\
    er/sie/es & \textbf{entschied}\\
    wir & \textbf{entschieden}\\
    ihr & \textbf{entschiedet}\\
    sie/Sie & \textbf{entschieden}\\
  \end{tabular}
  \hfill
  \begin{tabular}{rl}
     \multicolumn{2}{l}{\textbf{IMPERATIV} }\\
   & \textbf{entscheid(e) (du)!}\\
   & \textbf{entscheiden wir!}\\
   & \textbf{entscheidet (ihr)!}\\
   & \textbf{entscheiden Sie!}\\
   & \vspace{10pt}
  \end{tabular}

    \vspace{10pt}

 \begin{tabular}{rl}
     \multicolumn{2}{l}{\textbf{PERFEKT}}\\
  &  haben + \textbf{entschieden}\\
  \end{tabular}
\hfill
   \begin{tabular}{rl}
   
    \multicolumn{2}{l}{\textbf{FUTURE I} }\\
  &  werden + \textbf{entscheiden}\\
 
  \end{tabular}
\hfill
   \begin{tabular}{rl}
      \multicolumn{2}{l}{\textbf{PLUSQUAMPERFEKT}}\\
  &  hatten + \textbf{entschieden}\\
  \end{tabular}
  
   \vspace{10pt}

\begin{tabular}{lll}
\multicolumn{3}{l}{\textbf{MIT PRÄPOSITIONEN}}\\
& \textbf{---} & \\
\end{tabular}
\hfill
\begin{tabular}{lll}
\multicolumn{3}{l}{\textbf{TRENNBARE VERBEN}}\\
& \textbf{---} & \\
\end{tabular}

  \vspace{-5pt}
\end{verbentry}

% entsorgen (to dispose of)
\begin{verbentry}{
  style = {default},
  toc = {entsorgen (to dispose of)},
  name = {entsorgen},
  english = {to dispose of},
  type = {regular},
  auxiliary = {haben}
}
 
    
     \vspace{5pt}

    \hrule
    
    \vspace{5pt}

  \begin{tabular}{rl}
     \multicolumn{2}{l}{\textbf{PRÄSENS} }\\
    ich & \textbf{entsorge}\\
    du & \textbf{entsorgst}\\
    er/sie/es & \textbf{entsorgt}\\
    wir & \textbf{entsorgen}\\
    ihr & \textbf{entsorgt}\\
    sie/Sie & \textbf{entsorgen}\\
  \end{tabular}
  \hfill
     \begin{tabular}{rl}
    \multicolumn{2}{l}{\textbf{PRÄTERITUM} }\\
    ich & \textbf{entsorgte}\\
    du & \textbf{entsorgtest}\\
    er/sie/es & \textbf{entsorgte}\\
    wir & \textbf{entsorgten}\\
    ihr & \textbf{entsorgtet}\\
    sie/Sie & \textbf{entsorgten}\\
  \end{tabular}
  \hfill
  \begin{tabular}{rl}
     \multicolumn{2}{l}{\textbf{IMPERATIV} }\\
   & \textbf{entsorg(e) (du)!}\\
   & \textbf{entsorgen wir!}\\
   & \textbf{entsorgt (ihr)!}\\
   & \textbf{entsorgen Sie!}\\
   & \vspace{10pt}
  \end{tabular}

    \vspace{10pt}

 \begin{tabular}{rl}
     \multicolumn{2}{l}{\textbf{PERFEKT}}\\
   & haben + \textbf{entsorgt}\\
  \end{tabular}
\hfill
   \begin{tabular}{rl}
    \multicolumn{2}{l}{\textbf{FUTURE I} }\\
   & werde + \textbf{entsorgen}\\
  \end{tabular}
\hfill
   \begin{tabular}{rl}
      \multicolumn{2}{l}{\textbf{PLUSQUAMPERFEKT}}\\
   & hatten + \textbf{entsorgt}\\
  \end{tabular}

   \vspace{10pt}

   \begin{tabular}{lll}
      \multicolumn{3}{l}{\textbf{MIT PRÄPOSITIONEN}}\\
& \textbf{entsorgen über} + Akk & (to dispose of via)\\
& \textbf{entsorgen in} + Dat & (to dispose of in)\\
\end{tabular}
  \hfill
   \begin{tabular}{lll}
\multicolumn{3}{l}{\textbf{TRENNBARE VERBEN}}\\
& \textbf{---} & \\
\end{tabular}

  \vspace{-5pt}
\end{verbentry}

% entstehen (to arise)
\begin{verbentry}{
  style = {default},
  toc = {entstehen (to arise)},
  name = {entstehen},
  english = {to arise},
  type = {irregular},
  auxiliary = {sein}
}
 
    
     \vspace{5pt}

    \hrule
    
    \vspace{5pt}

  \begin{tabular}{rl}
     \multicolumn{2}{l}{\textbf{PRÄSENS} }\\
    ich & \textbf{entstehe}\\
    du & \textbf{entstehst}\\
    er/sie/es & \textbf{entsteht}\\
    wir & \textbf{entstehen}\\
    ihr & \textbf{entsteht}\\
    sie/Sie & \textbf{entstehen}\\
  \end{tabular}
  \hfill
     \begin{tabular}{rl}
    \multicolumn{2}{l}{\textbf{PRÄTERITUM} }\\
    ich & \textbf{entstand}\\
    du & \textbf{entstandest}\\
    er/sie/es & \textbf{entstand}\\
    wir & \textbf{entstanden}\\
    ihr & \textbf{entstandet}\\
    sie/Sie & \textbf{entstanden}\\
  \end{tabular}
  \hfill
  \begin{tabular}{rl}
     \multicolumn{2}{l}{\textbf{IMPERATIV} }\\
   & \textbf{entsteh(e) (du)!}\\
   & \textbf{entstehen wir!}\\
   & \textbf{entsteht (ihr)!}\\
   & \textbf{entstehen Sie!}\\
   & \vspace{10pt}
  \end{tabular}

    \vspace{10pt}

 \begin{tabular}{ll}
     \multicolumn{2}{l}{\textbf{PERFEKT}}\\
   & sein + \textbf{entstanden}\\
  \end{tabular}
\hfill
   \begin{tabular}{rl}
    \multicolumn{2}{l}{\textbf{FUTURE I} }\\
   & werden + \textbf{entstehen}\\
 
  \end{tabular}
\hfill
   \begin{tabular}{rl}
      \multicolumn{2}{l}{\textbf{PLUSQUAMPERFEKT}}\\
   & waren + \textbf{entstanden}\\
  \end{tabular}
  
\vspace{10pt}

\begin{tabular}{lll}
\multicolumn{3}{l}{\textbf{MIT PRÄPOSITIONEN}}\\
& \textbf{---} & \\
\end{tabular}
\hfill
\begin{tabular}{lll}
\multicolumn{3}{l}{\textbf{TRENNBARE VERBEN}}\\
& \textbf{---} & \\
\end{tabular}

  \vspace{-5pt}
\end{verbentry}

% entzünden (to ignite / inflame)
\begin{verbentry}{
  style = {default},
  toc = {entzünden (to ignite / inflame)},
  name = {entzünden},
  english = {to ignite / inflame},
  type = {regular},
  auxiliary = {haben}
}


\vspace{5pt}
\hrule
\vspace{5pt}

\begin{tabular}{rl}
\multicolumn{2}{l}{\textbf{PRÄSENS}}\\
ich & \textbf{entzünde}\\
du & \textbf{entzündest}\\
er/sie/es & \textbf{entzündet}\\
wir & \textbf{entzünden}\\
ihr & \textbf{entzündet}\\
sie/Sie & \textbf{entzünden}\\
\end{tabular}
\hfill
\begin{tabular}{rl}
\multicolumn{2}{l}{\textbf{PRÄTERITUM}}\\
ich & \textbf{entzündete}\\
du & \textbf{entzündetest}\\
er/sie/es & \textbf{entzündete}\\
wir & \textbf{entzündeten}\\
ihr & \textbf{entzündetet}\\
sie/Sie & \textbf{entzündeten}\\
\end{tabular}
\hfill
\begin{tabular}{rl}
\multicolumn{2}{l}{\textbf{IMPERATIV}}\\
& \textbf{entzünde (du)!}\\
& \textbf{entzünden wir!}\\
& \textbf{entzündet (ihr)!}\\
& \textbf{entzünden Sie!}\\
\end{tabular}

\vspace{10pt}

\begin{tabular}{rl}
\multicolumn{2}{l}{\textbf{PERFEKT}}\\
& haben + \textbf{entzündet}\\
\end{tabular}
\hfill
\begin{tabular}{rl}
\multicolumn{2}{l}{\textbf{FUTURE I}}\\
& werden + \textbf{entzünden}\\
\end{tabular}
\hfill
\begin{tabular}{rl}
\multicolumn{2}{l}{\textbf{PLUSQUAMPERFEKT}}\\
& hatten + \textbf{entzündet}\\
\end{tabular}

\vspace{10pt}

\begin{tabular}{lll}
\multicolumn{3}{l}{\textbf{MIT PRÄPOSITIONEN}}\\
& \textbf{entzünden an} + Dat & (to ignite on / by)\\
\end{tabular}
\hfill
\begin{tabular}{lll}
\multicolumn{3}{l}{\textbf{TRENNBARE VERBEN}}\\
& \textbf{---} & \\
\end{tabular}
\vspace{-5pt}
\end{verbentry}

% erfahren (to experience)
\begin{verbentry}{
  style = {acc},
  toc = {erfahren (to experience)},
  name = {erfahren},
  english = {to experience},
  type = {irregular, + Akkusativ},
  auxiliary = {haben}
}
 
    
     \vspace{5pt}

    \hrule
    
    \vspace{5pt}

  \begin{tabular}{rl}
     \multicolumn{2}{l}{\textbf{PRÄSENS} }\\
    ich & \textbf{erfahre}\\
    du & \textbf{erf\"ahrst}\\
    er/sie/es & \textbf{erf\"ahrt}\\
    wir & \textbf{erfahren}\\
    ihr & \textbf{erfahrt}\\
    sie/Sie & \textbf{erfahren}\\
  \end{tabular}
  \hfill
     \begin{tabular}{rl}
    \multicolumn{2}{l}{\textbf{PRÄTERITUM} }\\
    ich & \textbf{erfuhr}\\
    du & \textbf{erfuhrst}\\
    er/sie/es & \textbf{erfuhr}\\
    wir & \textbf{erfuhren}\\
    ihr & \textbf{erfuhrt}\\
    sie/Sie & \textbf{erfuhren}\\
  \end{tabular}
  \hfill
  \begin{tabular}{rl}
     \multicolumn{2}{l}{\textbf{IMPERATIV} }\\
   & \textbf{erfahr(e) (du)!}\\
   & \textbf{erfahren wir!}\\
   & \textbf{erfahrt (ihr)!}\\
   & \textbf{erfahren Sie!}\\
   & \vspace{10pt}
  \end{tabular}

    \vspace{10pt}

 \begin{tabular}{rl}
     \multicolumn{2}{l}{\textbf{PERFEKT}}\\
  &  haben + \textbf{erfahren}\\
  \end{tabular}
\hfill
   \begin{tabular}{rl}
   
    \multicolumn{2}{l}{\textbf{FUTURE I} }\\
  &  werden + \textbf{erfahren}\\
 
  \end{tabular}
\hfill
   \begin{tabular}{rl}
      \multicolumn{2}{l}{\textbf{PLUSQUAMPERFEKT}}\\
  &  hatten + \textbf{erfahren}\\
  \end{tabular}
  
  
\vspace{10pt}

\begin{tabular}{lll}
\multicolumn{3}{l}{\textbf{MIT PRÄPOSITIONEN}}\\
& \textbf{---} & \\
\end{tabular}
\hfill
\begin{tabular}{lll}
\multicolumn{3}{l}{\textbf{TRENNBARE VERBEN}}\\
& \textbf{---} & \\
\end{tabular}
\vspace{-5pt}
\end{verbentry}

% ergänzen (to complete / supplement)
\begin{verbentry}{
  style = {default},
  toc = {ergänzen (to complete / supplement)},
  name = {ergänzen},
  english = {to complete / supplement},
  type = {regular},
  auxiliary = {haben}
}
 

\vspace{5pt}
\hrule
\vspace{5pt}

\begin{tabular}{rl}
\multicolumn{2}{l}{\textbf{PRÄSENS}}\\
ich & \textbf{ergänze}\\
du & \textbf{ergänzt}\\
er/sie/es & \textbf{ergänzt}\\
wir & \textbf{ergänzen}\\
ihr & \textbf{ergänzt}\\
sie/Sie & \textbf{ergänzen}\\
\end{tabular}
\hfill
\begin{tabular}{rl}
\multicolumn{2}{l}{\textbf{PRÄTERITUM}}\\
ich & \textbf{ergänzte}\\
du & \textbf{ergänztest}\\
er/sie/es & \textbf{ergänzte}\\
wir & \textbf{ergänzten}\\
ihr & \textbf{ergänztet}\\
sie/Sie & \textbf{ergänzten}\\
\end{tabular}
\hfill
\begin{tabular}{rl}
\multicolumn{2}{l}{\textbf{IMPERATIV}}\\
& \textbf{ergänze (du)!}\\
& \textbf{ergänzen wir!}\\
& \textbf{ergänzt (ihr)!}\\
& \textbf{ergänzen Sie!}\\
\end{tabular}

\vspace{10pt}

\begin{tabular}{rl}
\multicolumn{2}{l}{\textbf{PERFEKT}}\\
& haben + \textbf{ergänzt}\\
\end{tabular}
\hfill
\begin{tabular}{rl}
\multicolumn{2}{l}{\textbf{FUTURE I}}\\
& werden + \textbf{ergänzen}\\
\end{tabular}
\hfill
\begin{tabular}{rl}
\multicolumn{2}{l}{\textbf{PLUSQUAMPERFEKT}}\\
& hatten + \textbf{ergänzt}\\
\end{tabular}

\vspace{10pt}

\begin{tabular}{lll}
\multicolumn{3}{l}{\textbf{MIT PRÄPOSITIONEN}}\\
& \textbf{ergänzen} + Akk & (to complete something)\\
\end{tabular}
\hfill
\begin{tabular}{lll}
\multicolumn{3}{l}{\textbf{TRENNBARE VERBEN}}\\
& \textbf{---} & \\
\end{tabular}

\vspace{-5pt}
\end{verbentry}

% erhalten (to receive / to preserve)
\begin{verbentry}{
  style = {default},
  toc = {erhalten (to receive / to preserve)},
  name = {erhalten},
  english = {to receive / to preserve},
  type = {irregular},
  auxiliary = {haben}
}
 
    
     \vspace{5pt}

    \hrule
    
    \vspace{5pt}

  \begin{tabular}{rl}
     \multicolumn{2}{l}{\textbf{PRÄSENS} }\\
    ich & \textbf{erhalte}\\
    du & \textbf{erhältst}\\
    er/sie/es & \textbf{erhält}\\
    wir & \textbf{erhalten}\\
    ihr & \textbf{erhaltet}\\
    sie/Sie & \textbf{erhalten}\\
  \end{tabular}
  \hfill
     \begin{tabular}{rl}
    \multicolumn{2}{l}{\textbf{PRÄTERITUM} }\\
    ich & \textbf{erhielt}\\
    du & \textbf{erhieltest}\\
    er/sie/es & \textbf{erhielt}\\
    wir & \textbf{erhielten}\\
    ihr & \textbf{erhieltet}\\
    sie/Sie & \textbf{erhielten}\\
  \end{tabular}
  \hfill
  \begin{tabular}{rl}
     \multicolumn{2}{l}{\textbf{IMPERATIV} }\\
   & \textbf{erhalte (du)!}\\
   & \textbf{erhalten wir!}\\
   & \textbf{erhaltet (ihr)!}\\
   & \textbf{erhalten Sie!}\\
   & \vspace{10pt}
  \end{tabular}

    \vspace{10pt}

 \begin{tabular}{rl}
     \multicolumn{2}{l}{\textbf{PERFEKT}}\\
   & haben + \textbf{erhalten}\\
  \end{tabular}
\hfill
   \begin{tabular}{rl}
    \multicolumn{2}{l}{\textbf{FUTURE I} }\\
   & werde + \textbf{erhalten}\\
  \end{tabular}
\hfill
   \begin{tabular}{rl}
      \multicolumn{2}{l}{\textbf{PLUSQUAMPERFEKT}}\\
   & hatten + \textbf{erhalten}\\
  \end{tabular}

   \vspace{10pt}

   \begin{tabular}{lll}
      \multicolumn{3}{l}{\textbf{MIT PRÄPOSITIONEN}}\\
& \textbf{erhalten von} + Dat & (to receive from)\\
& \textbf{erhalten für} + Akk & (to receive for)\\
\end{tabular}
  \hfill
   \begin{tabular}{lll}
\multicolumn{3}{l}{\textbf{TRENNBARE VERBEN}}\\
& \textbf{---} & \\
\end{tabular}

  \vspace{-5pt}
\end{verbentry}

% erinnern (to remember)
\begin{verbentry}{
  style = {acc},
  toc = {erinnern (to remember)},
  name = {erinnern},
  english = {to remember},
  type = {irregular, + Akkusativ},
  auxiliary = {haben}
}
 
    
     \vspace{5pt}

    \hrule
    
    \vspace{5pt}

  \begin{tabular}{rl}
     \multicolumn{2}{l}{\textbf{PRÄSENS} }\\
    ich & \textbf{erinnere}\\
    du & \textbf{erinnerst}\\
    er/sie/es & \textbf{erinnert}\\
    wir & \textbf{erinnern}\\
    ihr & \textbf{erinnert}\\
    sie/Sie & \textbf{erinnern}\\
  \end{tabular}
  \hfill
     \begin{tabular}{rl}
    \multicolumn{2}{l}{\textbf{PRÄTERITUM} }\\
    ich & \textbf{erinnerte}\\
    du & \textbf{erinnertest}\\
    er/sie/es & \textbf{erinnerte}\\
    wir & \textbf{erinnerten}\\
    ihr & \textbf{erinnertet}\\
    sie/Sie & \textbf{erinnerten}\\
  \end{tabular}
  \hfill
  \begin{tabular}{rl}
     \multicolumn{2}{l}{\textbf{IMPERATIV} }\\
   & \textbf{erinner(e) (du)!}\\
   & \textbf{erinnern wir!}\\
   & \textbf{erinnert (ihr)!}\\
   & \textbf{erinnern Sie!}\\
   & \vspace{10pt}
  \end{tabular}

    \vspace{10pt}

 \begin{tabular}{rl}
     \multicolumn{2}{l}{\textbf{PERFEKT}}\\
   & haben + \textbf{erinnert}\\
  \end{tabular}
\hfill
   \begin{tabular}{rl}
   
    \multicolumn{2}{l}{\textbf{FUTURE I} }\\
   & werden + \textbf{erinnern}\\
 
  \end{tabular}
\hfill
   \begin{tabular}{rl}
      \multicolumn{2}{l}{\textbf{PLUSQUAMPERFEKT}}\\
   & hatten + \textbf{erinnert}\\
  \end{tabular}
  
  \vspace{10pt}

\begin{tabular}{lll}
\multicolumn{3}{l}{\textbf{MIT PRÄPOSITIONEN}}\\
& \textbf{---} & \\
\end{tabular}
\hfill
\begin{tabular}{lll}
\multicolumn{3}{l}{\textbf{TRENNBARE VERBEN}}\\
& \textbf{---} & \\
\end{tabular}

  \vspace{-5pt}
\end{verbentry}

% erklären (to explain)
\begin{verbentry}{
  style = {dat},
  toc = {erklären (to explain)},
  name = {erklären},
  english = {to explain},
  type = {irregular, + Dativ},
  auxiliary = {haben}
}
 
    
     \vspace{5pt}

    \hrule
    
    \vspace{5pt}

  \begin{tabular}{rl}
     \multicolumn{2}{l}{\textbf{PRÄSENS} }\\
    ich & \textbf{erkläre}\\
    du & \textbf{erklärst}\\
    er/sie/es & \textbf{erklärt}\\
    wir & \textbf{erklären}\\
    ihr & \textbf{erklärt}\\
    sie/Sie & \textbf{erklären}\\
  \end{tabular}
  \hfill
     \begin{tabular}{rl}
    \multicolumn{2}{l}{\textbf{PRÄTERITUM} }\\
    ich & \textbf{erklärte}\\
    du & \textbf{erklärtest}\\
    er/sie/es & \textbf{erklärte}\\
    wir & \textbf{erklärten}\\
    ihr & \textbf{erklärtet}\\
    sie/Sie & \textbf{erklärten}\\
  \end{tabular}
  \hfill
  \begin{tabular}{rl}
     \multicolumn{2}{l}{\textbf{IMPERATIV} }\\
   & \textbf{erklär(e) (du)!}\\
   & \textbf{erklären wir!}\\
   & \textbf{erklärt (ihr)!}\\
   & \textbf{erklären Sie!}\\
   & \vspace{10pt}
  \end{tabular}

    \vspace{10pt}

 \begin{tabular}{rl}
     \multicolumn{2}{l}{\textbf{PERFEKT}}\\
  &  haben + \textbf{erklärt}\\
  \end{tabular}
\hfill
   \begin{tabular}{rl}
   
    \multicolumn{2}{l}{\textbf{FUTURE I} }\\
  &  werden + \textbf{erklären}\\
 
  \end{tabular}
\hfill
   \begin{tabular}{rl}
      \multicolumn{2}{l}{\textbf{PLUSQUAMPERFEKT}}\\
  &  hatten + \textbf{erklärt}\\
  \end{tabular}
  
  
\vspace{10pt}

\begin{tabular}{lll}
\multicolumn{3}{l}{\textbf{MIT PRÄPOSITIONEN}}\\
& \textbf{---} & \\
\end{tabular}
\hfill
\begin{tabular}{lll}
\multicolumn{3}{l}{\textbf{TRENNBARE VERBEN}}\\
& \textbf{---} & \\
\end{tabular}
\vspace{-5pt}
\end{verbentry}

% erkunden (to explore)
\begin{verbentry}{
  style = {acc},
  toc = {erkunden (to explore)},
  name = {erkunden},
  english = {to explore},
  type = {regular, + Akkusativ},
  auxiliary = {haben}
}


\vspace{5pt}
\hrule
\vspace{5pt}

\begin{tabular}{rl}
\multicolumn{2}{l}{\textbf{PRÄSENS}}\\
ich & \textbf{erkunde}\\
du & \textbf{erkundest}\\
er/sie/es & \textbf{erkundet}\\
wir & \textbf{erkunden}\\
ihr & \textbf{erkundet}\\
sie/Sie & \textbf{erkunden}\\
\end{tabular}
\hfill
\begin{tabular}{rl}
\multicolumn{2}{l}{\textbf{PRÄTERITUM}}\\
ich & \textbf{erkundete}\\
du & \textbf{erkundetest}\\
er/sie/es & \textbf{erkundete}\\
wir & \textbf{erkundeten}\\
ihr & \textbf{erkundetet}\\
sie/Sie & \textbf{erkundeten}\\
\end{tabular}
\hfill
\begin{tabular}{rl}
\multicolumn{2}{l}{\textbf{IMPERATIV}}\\
& \textbf{erkunde (du)!}\\
& \textbf{erkunden wir!}\\
& \textbf{erkundet (ihr)!}\\
& \textbf{erkunden Sie!}\\
\end{tabular}

\vspace{10pt}

\begin{tabular}{rl}
\multicolumn{2}{l}{\textbf{PERFEKT}}\\
& haben + \textbf{erkundet}\\
\end{tabular}
\hfill
\begin{tabular}{rl}
\multicolumn{2}{l}{\textbf{FUTURE I}}\\
& werden + \textbf{erkunden}\\
\end{tabular}
\hfill
\begin{tabular}{rl}
\multicolumn{2}{l}{\textbf{PLUSQUAMPERFEKT}}\\
& hatten + \textbf{erkundet}\\
\end{tabular}

\vspace{10pt}

\begin{tabular}{lll}
\multicolumn{3}{l}{\textbf{MIT PRÄPOSITIONEN}}\\
& \textbf{---} & \\
\end{tabular}
\hfill
\begin{tabular}{lll}
\multicolumn{3}{l}{\textbf{TRENNBARE VERBEN}}\\
& \textbf{---} & \\
\end{tabular}

\vspace{-5pt}
\end{verbentry}

% erlauben (to allow)
\begin{verbentry}{
  style = {acc},
  toc = {erlauben (to allow)},
  name = {erlauben},
  english = {to allow},
  type = {regular, + Akkusativ},
  auxiliary = {haben}
}
 
    
     \vspace{5pt}

    \hrule
    
    \vspace{5pt}

  \begin{tabular}{rl}
     \multicolumn{2}{l}{\textbf{PRÄSENS} }\\
    ich & \textbf{erlaube}\\
    du & \textbf{erlaubst}\\
    er/sie/es & \textbf{erlaubt}\\
    wir & \textbf{erlauben}\\
    ihr & \textbf{erlaubt}\\
    sie/Sie & \textbf{erlauben}\\
  \end{tabular}
  \hfill
     \begin{tabular}{rl}
    \multicolumn{2}{l}{\textbf{PRÄTERITUM} }\\
    ich & \textbf{erlaubte}\\
    du & \textbf{erlaubtest}\\
    er/sie/es & \textbf{erlaubte}\\
    wir & \textbf{erlaubten}\\
    ihr & \textbf{erlaubtet}\\
    sie/Sie & \textbf{erlaubten}\\
  \end{tabular}
  \hfill
  \begin{tabular}{rl}
     \multicolumn{2}{l}{\textbf{IMPERATIV} }\\
   & \textbf{erlaub(e) (du)!}\\
   & \textbf{erlauben wir!}\\
   & \textbf{erlaubt (ihr)!}\\
   & \textbf{erlauben Sie!}\\
   & \vspace{10pt}
  \end{tabular}

    \vspace{10pt}

 \begin{tabular}{rl}
     \multicolumn{2}{l}{\textbf{PERFEKT}}\\
   & haben + \textbf{erlaubt}\\
  \end{tabular}
\hfill
   \begin{tabular}{rl}
   
    \multicolumn{2}{l}{\textbf{FUTURE I} }\\
  &  werden + \textbf{erlauben}\\
 
  \end{tabular}
\hfill
   \begin{tabular}{rl}
      \multicolumn{2}{l}{\textbf{PLUSQUAMPERFEKT}}\\
  &  hatten + \textbf{erlaubt}\\
  \end{tabular}
  
  \vspace{10pt}

\begin{tabular}{lll}
\multicolumn{3}{l}{\textbf{MIT PRÄPOSITIONEN}}\\
& \textbf{---} & \\
\end{tabular}
\hfill
\begin{tabular}{lll}
\multicolumn{3}{l}{\textbf{TRENNBARE VERBEN}}\\
& \textbf{---} & \\
\end{tabular}

  \vspace{-5pt}
\end{verbentry}

% erleben (to experience)
\begin{verbentry}{
  style = {acc},
  toc = {erleben (to experience)},
  name = {erleben},
  english = {to experience},
  type = {regular, + Akkusativ},
  auxiliary = {haben}
}


\vspace{5pt}
\hrule
\vspace{5pt}

\begin{tabular}{rl}
\multicolumn{2}{l}{\textbf{PRÄSENS}}\\
ich & \textbf{erlebe}\\
du & \textbf{erlebst}\\
er/sie/es & \textbf{erlebt}\\
wir & \textbf{erleben}\\
ihr & \textbf{erlebt}\\
sie/Sie & \textbf{erleben}\\
\end{tabular}
\hfill
\begin{tabular}{rl}
\multicolumn{2}{l}{\textbf{PRÄTERITUM}}\\
ich & \textbf{erlebte}\\
du & \textbf{erlebtest}\\
er/sie/es & \textbf{erlebte}\\
wir & \textbf{erlebten}\\
ihr & \textbf{erlebtet}\\
sie/Sie & \textbf{erlebten}\\
\end{tabular}
\hfill
\begin{tabular}{rl}
\multicolumn{2}{l}{\textbf{IMPERATIV}}\\
& \textbf{erleb(e) (du)!}\\
& \textbf{erleben wir!}\\
& \textbf{erlebt (ihr)!}\\
& \textbf{erleben Sie!}\\
\end{tabular}

\vspace{10pt}

\begin{tabular}{rl}
\multicolumn{2}{l}{\textbf{PERFEKT}}\\
& haben + \textbf{erlebt}\\
\end{tabular}
\hfill
\begin{tabular}{rl}
\multicolumn{2}{l}{\textbf{FUTURE I}}\\
& werden + \textbf{erleben}\\
\end{tabular}
\hfill
\begin{tabular}{rl}
\multicolumn{2}{l}{\textbf{PLUSQUAMPERFEKT}}\\
& hatten + \textbf{erlebt}\\
\end{tabular}

\vspace{10pt}

\begin{tabular}{lll}
\multicolumn{3}{l}{\textbf{MIT PRÄPOSITIONEN}}\\
& \textbf{---} & \\
\end{tabular}
\hfill
\begin{tabular}{lll}
\multicolumn{3}{l}{\textbf{TRENNBARE VERBEN}}\\
& \textbf{---} & \\
\end{tabular}

\vspace{-5pt}
\end{verbentry}

% erledigen (to complete / take care of)
\begin{verbentry}{
  style = {default},
  toc = {erledigen (to complete / take care of)},
  name = {erledigen},
  english = {to complete / take care of},
  type = {regular},
  auxiliary = {haben}
}
 

\vspace{5pt}
\hrule
\vspace{5pt}

\begin{tabular}{rl}
\multicolumn{2}{l}{\textbf{PRÄSENS}}\\
ich & \textbf{erledige}\\
du & \textbf{erledigst}\\
er/sie/es & \textbf{erledigt}\\
wir & \textbf{erledigen}\\
ihr & \textbf{erledigt}\\
sie/Sie & \textbf{erledigen}\\
\end{tabular}
\hfill
\begin{tabular}{rl}
\multicolumn{2}{l}{\textbf{PRÄTERITUM}}\\
ich & \textbf{erledigte}\\
du & \textbf{erledigtest}\\
er/sie/es & \textbf{erledigte}\\
wir & \textbf{erledigten}\\
ihr & \textbf{erledigtet}\\
sie/Sie & \textbf{erledigten}\\
\end{tabular}
\hfill
\begin{tabular}{rl}
\multicolumn{2}{l}{\textbf{IMPERATIV}}\\
& \textbf{erledige (du)!}\\
& \textbf{erledigen wir!}\\
& \textbf{erledigt (ihr)!}\\
& \textbf{erledigen Sie!}\\
\end{tabular}

\vspace{10pt}

\begin{tabular}{rl}
\multicolumn{2}{l}{\textbf{PERFEKT}}\\
& haben + \textbf{erledigt}\\
\end{tabular}
\hfill
\begin{tabular}{rl}
\multicolumn{2}{l}{\textbf{FUTURE I}}\\
& werden + \textbf{erledigen}\\
\end{tabular}
\hfill
\begin{tabular}{rl}
\multicolumn{2}{l}{\textbf{PLUSQUAMPERFEKT}}\\
& hatten + \textbf{erledigt}\\
\end{tabular}

\vspace{10pt}

\begin{tabular}{lll}
\multicolumn{3}{l}{\textbf{MIT PRÄPOSITIONEN}}\\
& \textbf{erledigen} + Akk & (to complete something)\\
\end{tabular}
\hfill
\begin{tabular}{lll}
\multicolumn{3}{l}{\textbf{TRENNBARE VERBEN}}\\
& \textbf{---} & \\
\end{tabular}

\vspace{-5pt}
\end{verbentry}

% erlernen (to learn)
\begin{verbentry}{
  style = {acc},
  toc = {erlernen (to learn)},
  name = {erlernen},
  english = {to learn / acquire a skill},
  type = {regular, + Akkusativ},
  auxiliary = {haben}
}


\vspace{5pt}
\hrule
\vspace{5pt}

\begin{tabular}{rl}
\multicolumn{2}{l}{\textbf{PRÄSENS}}\\
ich & \textbf{erlerne}\\
du & \textbf{erlernst}\\
er/sie/es & \textbf{erlernt}\\
wir & \textbf{erlernen}\\
ihr & \textbf{erlernt}\\
sie/Sie & \textbf{erlernen}\\
\end{tabular}
\hfill
\begin{tabular}{rl}
\multicolumn{2}{l}{\textbf{PRÄTERITUM}}\\
ich & \textbf{erlernte}\\
du & \textbf{erlerntest}\\
er/sie/es & \textbf{erlernte}\\
wir & \textbf{erlernten}\\
ihr & \textbf{erlerntet}\\
sie/Sie & \textbf{erlernten}\\
\end{tabular}
\hfill
\begin{tabular}{rl}
\multicolumn{2}{l}{\textbf{IMPERATIV}}\\
& \textbf{erlern(e) (du)!}\\
& \textbf{erlernen wir!}\\
& \textbf{erlernt (ihr)!}\\
& \textbf{erlernen Sie!}\\
\end{tabular}

\vspace{10pt}

\begin{tabular}{rl}
\multicolumn{2}{l}{\textbf{PERFEKT}}\\
& haben + \textbf{erlernt}\\
\end{tabular}
\hfill
\begin{tabular}{rl}
\multicolumn{2}{l}{\textbf{FUTURE I}}\\
& werden + \textbf{erlernen}\\
\end{tabular}
\hfill
\begin{tabular}{rl}
\multicolumn{2}{l}{\textbf{PLUSQUAMPERFEKT}}\\
& hatten + \textbf{erlernt}\\
\end{tabular}

\vspace{10pt}

\begin{tabular}{lll}
\multicolumn{3}{l}{\textbf{MIT PRÄPOSITIONEN}}\\
& \textbf{---} & \\
\end{tabular}
\hfill
\begin{tabular}{lll}
\multicolumn{3}{l}{\textbf{TRENNBARE VERBEN}}\\
& \textbf{---} & \\
\end{tabular}

\vspace{-5pt}
\end{verbentry}

% eröffnen (to open / inaugurate)
\begin{verbentry}{
  style = {default},
  toc = {eröffnen (to open / inaugurate)},
  name = {eröffnen},
  english = {to open / inaugurate},
  type = {regular},
  auxiliary = {haben}
}


\vspace{5pt}
\hrule
\vspace{5pt}

\begin{tabular}{rl}
\multicolumn{2}{l}{\textbf{PRÄSENS}}\\
ich & \textbf{eröffne}\\
du & \textbf{eröffnest}\\
er/sie/es & \textbf{eröffnet}\\
wir & \textbf{eröffnen}\\
ihr & \textbf{eröffnet}\\
sie/Sie & \textbf{eröffnen}\\
\end{tabular}
\hfill
\begin{tabular}{rl}
\multicolumn{2}{l}{\textbf{PRÄTERITUM}}\\
ich & \textbf{eröffnete}\\
du & \textbf{eröffnetest}\\
er/sie/es & \textbf{eröffnete}\\
wir & \textbf{eröffneten}\\
ihr & \textbf{eröffnetet}\\
sie/Sie & \textbf{eröffneten}\\
\end{tabular}
\hfill
\begin{tabular}{rl}
\multicolumn{2}{l}{\textbf{IMPERATIV}}\\
& \textbf{eröffne (du)!}\\
& \textbf{eröffnen wir!}\\
& \textbf{eröffnet (ihr)!}\\
& \textbf{eröffnen Sie!}\\
\end{tabular}

\vspace{10pt}

\begin{tabular}{rl}
\multicolumn{2}{l}{\textbf{PERFEKT}}\\
& haben + \textbf{eröffnet}\\
\end{tabular}
\hfill
\begin{tabular}{rl}
\multicolumn{2}{l}{\textbf{FUTURE I}}\\
& werden + \textbf{eröffnen}\\
\end{tabular}
\hfill
\begin{tabular}{rl}
\multicolumn{2}{l}{\textbf{PLUSQUAMPERFEKT}}\\
& hatten + \textbf{eröffnet}\\
\end{tabular}

\vspace{10pt}

\begin{tabular}{lll}
\multicolumn{3}{l}{\textbf{MIT PRÄPOSITIONEN}}\\
& \textbf{eröffnen für} + Akk & (to open for)\\
& \textbf{eröffnen zu} + Dat & (to inaugurate at / on)\\
\end{tabular}
\hfill
\begin{tabular}{lll}
\multicolumn{3}{l}{\textbf{TRENNBARE VERBEN}}\\
& \textbf{---} & \\
\end{tabular}
\vspace{-5pt}
\end{verbentry}

% erreichen (to reach, achieve)
\begin{verbentry}{
  style = {acc},
  toc = {erreichen (to reach, achieve)},
  name = {erreichen},
  english = {to reach, achieve},
  type = {regular, + Akkusativ},
  auxiliary = {haben}
}
 
    
     \vspace{5pt}

    \hrule
    
    \vspace{5pt}

  \begin{tabular}{rl}
     \multicolumn{2}{l}{\textbf{PRÄSENS} }\\
    ich & \textbf{erreiche}\\
    du & \textbf{erreichst}\\
    er/sie/es & \textbf{erreicht}\\
    wir & \textbf{erreichen}\\
    ihr & \textbf{erreicht}\\
    sie/Sie & \textbf{erreichen}\\
  \end{tabular}
  \hfill
     \begin{tabular}{rl}
    \multicolumn{2}{l}{\textbf{PRÄTERITUM} }\\
    ich & \textbf{erreichte}\\
    du & \textbf{erreichtest}\\
    er/sie/es & \textbf{erreichte}\\
    wir & \textbf{erreichten}\\
    ihr & \textbf{erreichtet}\\
    sie/Sie & \textbf{erreichten}\\
  \end{tabular}
  \hfill
  \begin{tabular}{rl}
     \multicolumn{2}{l}{\textbf{IMPERATIV} }\\
   & \textbf{erreich(e) (du)!}\\
   & \textbf{erreichen wir!}\\
   & \textbf{erreicht (ihr)!}\\
   & \textbf{erreichen Sie!}\\
   & \vspace{10pt}
  \end{tabular}

    \vspace{10pt}

 \begin{tabular}{rl}
     \multicolumn{2}{l}{\textbf{PERFEKT}}\\
  &  haben + \textbf{erreicht}\\
  \end{tabular}
\hfill
   \begin{tabular}{rl}
   
    \multicolumn{2}{l}{\textbf{FUTURE I} }\\
  &  werden + \textbf{erreichen}\\
 
  \end{tabular}
\hfill
   \begin{tabular}{rl}
      \multicolumn{2}{l}{\textbf{PLUSQUAMPERFEKT}}\\
  &  hatten + \textbf{erreicht}\\
  \end{tabular}
  
  
\vspace{10pt}

\begin{tabular}{lll}
\multicolumn{3}{l}{\textbf{MIT PRÄPOSITIONEN}}\\
& \textbf{---} & \\
\end{tabular}
\hfill
\begin{tabular}{lll}
\multicolumn{3}{l}{\textbf{TRENNBARE VERBEN}}\\
& \textbf{---} & \\
\end{tabular}
\vspace{-5pt}
\end{verbentry}

% erneuern (to renew / renovate)
\begin{verbentry}{
  style = {default},
  toc = {erneuern (to renew / renovate)},
  name = {erneuern},
  english = {to renew / renovate},
  type = {regular},
  auxiliary = {haben}
}
 
    
     \vspace{5pt}

    \hrule
    
    \vspace{5pt}

  \begin{tabular}{rl}
     \multicolumn{2}{l}{\textbf{PRÄSENS} }\\
    ich & \textbf{erneuere}\\
    du & \textbf{erneuerst}\\
    er/sie/es & \textbf{erneuert}\\
    wir & \textbf{erneuern}\\
    ihr & \textbf{erneuert}\\
    sie/Sie & \textbf{erneuern}\\
  \end{tabular}
  \hfill
  \begin{tabular}{rl}
    \multicolumn{2}{l}{\textbf{PRÄTERITUM} }\\
    ich & \textbf{erneuerte}\\
    du & \textbf{erneuertest}\\
    er/sie/es & \textbf{erneuerte}\\
    wir & \textbf{erneuerten}\\
    ihr & \textbf{erneuertet}\\
    sie/Sie & \textbf{erneuerten}\\
  \end{tabular}
  \hfill
  \begin{tabular}{rl}
     \multicolumn{2}{l}{\textbf{IMPERATIV} }\\
   & \textbf{erneuere (du)!}\\
   & \textbf{erneuern wir!}\\
   & \textbf{erneuert (ihr)!}\\
   & \textbf{erneuern Sie!}\\
   & \vspace{10pt}
  \end{tabular}

 \vspace{10pt}

 \begin{tabular}{rl}
     \multicolumn{2}{l}{\textbf{PERFEKT}}\\
   & haben + \textbf{erneuert}\\
  \end{tabular}
\hfill
\begin{tabular}{rl}
    \multicolumn{2}{l}{\textbf{FUTURE I}}\\
   & werden + \textbf{erneuern}\\
  \end{tabular}
\hfill
\begin{tabular}{rl}
    \multicolumn{2}{l}{\textbf{PLUSQUAMPERFEKT}}\\
   & hatten + \textbf{erneuert}\\
  \end{tabular}

 \vspace{10pt}

\begin{tabular}{lll}
      \multicolumn{3}{l}{\textbf{MIT PRÄPOSITIONEN}}\\
& \textbf{erneuern von} + Dat & (to renew from / of)\\
& \textbf{erneuern durch} + Akk & (to renew through / by)\\
& \textbf{erneuern zu} + Dat & (to renew into)\\
\end{tabular}
  \hfill
\begin{tabular}{lll}
      \multicolumn{3}{l}{\textbf{TRENNBARE VERBEN}}\\
 & \textbf{ein$|$erneuern} & (to renew in / reinstall)\\
\vspace{20pt}
  \end{tabular}


\vspace{-5pt}
\end{verbentry}

% erscheinen (to appear)
\begin{verbentry}{
  style = {dat},
  toc = {erscheinen (to appear)},
  name = {erscheinen},
  english = {to appear},
  type = {irregular, inseparable},
  auxiliary = {sein}
}
 
    
     \vspace{5pt}

    \hrule
    
    \vspace{5pt}

  \begin{tabular}{rl}
     \multicolumn{2}{l}{\textbf{PRÄSENS} }\\
    ich & \textbf{erscheine}\\
    du & \textbf{erscheinst}\\
    er/sie/es & \textbf{erscheint}\\
    wir & \textbf{erscheinen}\\
    ihr & \textbf{erscheint}\\
    sie/Sie & \textbf{erscheinen}\\
  \end{tabular}
  \hfill
     \begin{tabular}{rl}
    \multicolumn{2}{l}{\textbf{PRÄTERITUM} }\\
    ich & \textbf{erschien}\\
    du & \textbf{erschienst}\\
    er/sie/es & \textbf{erschien}\\
    wir & \textbf{erschienen}\\
    ihr & \textbf{erschient}\\
    sie/Sie & \textbf{erschienen}\\
  \end{tabular}
  \hfill
  \begin{tabular}{rl}
     \multicolumn{2}{l}{\textbf{IMPERATIV} }\\
   & \textbf{erschein(e) (du)!}\\
   & \textbf{erscheinen wir!}\\
   & \textbf{erscheint (ihr)!}\\
   & \textbf{erscheinen Sie!}\\
   & \vspace{10pt}
  \end{tabular}

    \vspace{10pt}

 \begin{tabular}{rl}
     \multicolumn{2}{l}{\textbf{PERFEKT}}\\
   & sein + \textbf{erschienen}\\
  \end{tabular}
\hfill
   \begin{tabular}{rl}
   
    \multicolumn{2}{l}{\textbf{FUTURE I} }\\
   & werden + \textbf{erscheinen}\\
 
  \end{tabular}
\hfill
   \begin{tabular}{rl}
      \multicolumn{2}{l}{\textbf{PLUSQUAMPERFEKT}}\\
   & waren + \textbf{erschienen}\\
  \end{tabular}

   \vspace{10pt}

   \begin{tabular}{lll}
      \multicolumn{3}{l}{\textbf{MIT PRÄPOSITIONEN}}\\
& \textbf{erscheinen in} + Dat & (to appear in)\\
& \textbf{erscheinen vor} + Dat & (to appear before)\\
\end{tabular}
  \hfill
   \begin{tabular}{lll}
      \multicolumn{3}{l}{\textbf{TRENNBARE VERBEN}}\\
 & -- & --\\
\vspace{10pt}
  \end{tabular}


  \vspace{-5pt}
\end{verbentry}

% ersetzen (to replace)
\begin{verbentry}{
  style = {default},
  toc = {ersetzen (to replace)},
  name = {ersetzen},
  english = {to replace},
  type = {regular, inseparable},
  auxiliary = {haben}
}
 
    
     \vspace{5pt}

    \hrule
    
    \vspace{5pt}

  \begin{tabular}{rl}
     \multicolumn{2}{l}{\textbf{PRÄSENS} }\\
    ich & \textbf{ersetze}\\
    du & \textbf{ersetzt}\\
    er/sie/es & \textbf{ersetzt}\\
    wir & \textbf{ersetzen}\\
    ihr & \textbf{ersetzt}\\
    sie/Sie & \textbf{ersetzen}\\
  \end{tabular}
  \hfill
     \begin{tabular}{rl}
    \multicolumn{2}{l}{\textbf{PRÄTERITUM} }\\
    ich & \textbf{ersetzte}\\
    du & \textbf{ersetztest}\\
    er/sie/es & \textbf{ersetzte}\\
    wir & \textbf{ersetzten}\\
    ihr & \textbf{ersetztet}\\
    sie/Sie & \textbf{ersetzten}\\
  \end{tabular}
  \hfill
  \begin{tabular}{rl}
     \multicolumn{2}{l}{\textbf{IMPERATIV} }\\
   & \textbf{ersetz(e) (du)!}\\
   & \textbf{ersetzen wir!}\\
   & \textbf{ersetzt (ihr)!}\\
   & \textbf{ersetzen Sie!}\\
   & \vspace{10pt}
  \end{tabular}

    \vspace{10pt}

 \begin{tabular}{rl}
     \multicolumn{2}{l}{\textbf{PERFEKT}}\\
   & haben + \textbf{ersetzt}\\
  \end{tabular}
\hfill
   \begin{tabular}{rl}
    \multicolumn{2}{l}{\textbf{FUTURE I} }\\
   & werden + \textbf{ersetzen}\\
  \end{tabular}
\hfill
   \begin{tabular}{rl}
      \multicolumn{2}{l}{\textbf{PLUSQUAMPERFEKT}}\\
   & hatten + \textbf{ersetzt}\\
  \end{tabular}

   \vspace{10pt}

   \begin{tabular}{lll}
      \multicolumn{3}{l}{\textbf{MIT PRÄPOSITIONEN}}\\
& \textbf{ersetzen durch} + Akk & (to replace with / by)\\
& \textbf{ersetzen mit} + Dat & (to replace with)\\
\end{tabular}
  \hfill
   \begin{tabular}{lll}
\multicolumn{3}{l}{\textbf{TRENNBARE VERBEN}}\\
& \textbf{---} & \\
\end{tabular}

  \vspace{-5pt}
\end{verbentry}

% erstellen (to create / to prepare)
\begin{verbentry}{
  style = {default},
  toc = {erstellen (to create / to prepare)},
  name = {erstellen},
  english = {to create / to prepare},
  type = {regular, inseparable},
  auxiliary = {haben}
}
 
    
     \vspace{5pt}

    \hrule
    
    \vspace{5pt}

  \begin{tabular}{rl}
     \multicolumn{2}{l}{\textbf{PRÄSENS} }\\
    ich & \textbf{erstelle}\\
    du & \textbf{erstellst}\\
    er/sie/es & \textbf{erstellt}\\
    wir & \textbf{erstellen}\\
    ihr & \textbf{erstellt}\\
    sie/Sie & \textbf{erstellen}\\
  \end{tabular}
  \hfill
     \begin{tabular}{rl}
    \multicolumn{2}{l}{\textbf{PRÄTERITUM} }\\
    ich & \textbf{erstellte}\\
    du & \textbf{erstelltest}\\
    er/sie/es & \textbf{erstellte}\\
    wir & \textbf{erstellten}\\
    ihr & \textbf{erstelltet}\\
    sie/Sie & \textbf{erstellten}\\
  \end{tabular}
  \hfill
  \begin{tabular}{rl}
     \multicolumn{2}{l}{\textbf{IMPERATIV} }\\
   & \textbf{erstelle (du)!}\\
   & \textbf{erstellen wir!}\\
   & \textbf{erstellt (ihr)!}\\
   & \textbf{erstellen Sie!}\\
   & \vspace{10pt}
  \end{tabular}

    \vspace{10pt}

 \begin{tabular}{rl}
     \multicolumn{2}{l}{\textbf{PERFEKT}}\\
   & haben + \textbf{erstellt}\\
  \end{tabular}
\hfill
   \begin{tabular}{rl}
    \multicolumn{2}{l}{\textbf{FUTURE I} }\\
   & werden + \textbf{erstellen}\\
  \end{tabular}
\hfill
   \begin{tabular}{rl}
      \multicolumn{2}{l}{\textbf{PLUSQUAMPERFEKT}}\\
   & hatten + \textbf{erstellt}\\
  \end{tabular}

   \vspace{10pt}

   \begin{tabular}{lll}
      \multicolumn{3}{l}{\textbf{MIT PRÄPOSITIONEN}}\\
& \textbf{erstellen für} + Akk & (to create for)\\
& \textbf{erstellen mit} + Dat & (to create with)\\
\end{tabular}
  \hfill
   \begin{tabular}{lll}
\multicolumn{3}{l}{\textbf{TRENNBARE VERBEN}}\\
& \textbf{---} & \\
\end{tabular}

  \vspace{-5pt}
\end{verbentry}

% erwarten (to expect)
\begin{verbentry}{
  style = {default},
  toc = {erwarten (to expect)},
  name = {erwarten},
  english = {to expect},
  type = {regular},
  auxiliary = {haben}
}


\vspace{5pt}
\hrule
\vspace{5pt}

\begin{tabular}{rl}
\multicolumn{2}{l}{\textbf{PRÄSENS}}\\
ich & \textbf{erwarte}\\
du & \textbf{erwartest}\\
er/sie/es & \textbf{erwartet}\\
wir & \textbf{erwarten}\\
ihr & \textbf{erwartet}\\
sie/Sie & \textbf{erwarten}\\
\end{tabular}
\hfill
\begin{tabular}{rl}
\multicolumn{2}{l}{\textbf{PRÄTERITUM}}\\
ich & \textbf{erwartete}\\
du & \textbf{erwartetest}\\
er/sie/es & \textbf{erwartete}\\
wir & \textbf{erwarteten}\\
ihr & \textbf{erwartetet}\\
sie/Sie & \textbf{erwarteten}\\
\end{tabular}
\hfill
\begin{tabular}{rl}
\multicolumn{2}{l}{\textbf{IMPERATIV}}\\
& \textbf{erwarte (du)!}\\
& \textbf{erwarten wir!}\\
& \textbf{erwartet (ihr)!}\\
& \textbf{erwarten Sie!}\\
\end{tabular}

\vspace{10pt}

\begin{tabular}{rl}
\multicolumn{2}{l}{\textbf{PERFEKT}}\\
& haben + \textbf{erwartet}\\
\end{tabular}
\hfill
\begin{tabular}{rl}
\multicolumn{2}{l}{\textbf{FUTURE I}}\\
& werden + \textbf{erwarten}\\
\end{tabular}
\hfill
\begin{tabular}{rl}
\multicolumn{2}{l}{\textbf{PLUSQUAMPERFEKT}}\\
& hatten + \textbf{erwartet}\\
\end{tabular}

\vspace{10pt}

\begin{tabular}{lll}
\multicolumn{3}{l}{\textbf{MIT PRÄPOSITIONEN}}\\
& \textbf{erwarten von} + Dat & (to expect from)\\
\end{tabular}
\hfill
\begin{tabular}{lll}
\multicolumn{3}{l}{\textbf{TRENNBARE VERBEN}}\\
& \textbf{---} & \\
\end{tabular}
\vspace{-5pt}
\end{verbentry}

% erweitern (to expand)
\begin{verbentry}{
  style = {acc},
  toc = {erweitern (to expand)},
  name = {erweitern},
  english = {to expand},
  type = {regular, + Akkusativ},
  auxiliary = {haben}
}


\vspace{5pt}
\hrule
\vspace{5pt}

\begin{tabular}{rl}
\multicolumn{2}{l}{\textbf{PRÄSENS}}\\
ich & \textbf{erweitere}\\
du & \textbf{erweiterst}\\
er/sie/es & \textbf{erweitert}\\
wir & \textbf{erweitern}\\
ihr & \textbf{erweitert}\\
sie/Sie & \textbf{erweitern}\\
\end{tabular}
\hfill
\begin{tabular}{rl}
\multicolumn{2}{l}{\textbf{PRÄTERITUM}}\\
ich & \textbf{erweiterte}\\
du & \textbf{erweitertest}\\
er/sie/es & \textbf{erweiterte}\\
wir & \textbf{erweiterten}\\
ihr & \textbf{erweitertet}\\
sie/Sie & \textbf{erweiterten}\\
\end{tabular}
\hfill
\begin{tabular}{rl}
\multicolumn{2}{l}{\textbf{IMPERATIV}}\\
& \textbf{erweiter(e) (du)!}\\
& \textbf{erweitern wir!}\\
& \textbf{erweitert (ihr)!}\\
& \textbf{erweitern Sie!}\\
\end{tabular}

\vspace{10pt}

\begin{tabular}{rl}
\multicolumn{2}{l}{\textbf{PERFEKT}}\\
& haben + \textbf{erweitert}\\
\end{tabular}
\hfill
\begin{tabular}{rl}
\multicolumn{2}{l}{\textbf{FUTURE I}}\\
& werden + \textbf{erweitern}\\
\end{tabular}
\hfill
\begin{tabular}{rl}
\multicolumn{2}{l}{\textbf{PLUSQUAMPERFEKT}}\\
& hatten + \textbf{erweitert}\\
\end{tabular}

\vspace{10pt}

\begin{tabular}{lll}
\multicolumn{3}{l}{\textbf{MIT PRÄPOSITIONEN}}\\
& \textbf{erweitern um} + Akk & (to expand by)\\
\end{tabular}
\hfill
\begin{tabular}{lll}
\multicolumn{3}{l}{\textbf{TRENNBARE VERBEN}}\\
& \textbf{---} & \\
\end{tabular}

\vspace{-5pt}
\end{verbentry}

% erzählen (to tell)
\begin{verbentry}{
  style = {dat},
  toc = {erzählen (to tell)},
  name = {erzählen},
  english = {to tell},
  type = {regular, + Dativ},
  auxiliary = {haben}
}
 
    
     \vspace{5pt}

    \hrule
    
    \vspace{5pt}

  \begin{tabular}{rl}
     \multicolumn{2}{l}{\textbf{PRÄSENS} }\\
    ich & \textbf{erzähle}\\
    du & \textbf{erzählst}\\
    er/sie/es & \textbf{erzählt}\\
    wir & \textbf{erzählen}\\
    ihr & \textbf{erzählt}\\
    sie/Sie & \textbf{erzählen}\\
  \end{tabular}
  \hfill
     \begin{tabular}{rl}
    \multicolumn{2}{l}{\textbf{PRÄTERITUM} }\\
    ich & \textbf{erzählte}\\
    du & \textbf{erzähltest}\\
    er/sie/es & \textbf{erzählte}\\
    wir & \textbf{erzählten}\\
    ihr & \textbf{erzähltet}\\
    sie/Sie & \textbf{erzählten}\\
  \end{tabular}
  \hfill
  \begin{tabular}{rl}
     \multicolumn{2}{l}{\textbf{IMPERATIV} }\\
   & \textbf{erzähl(e) (du)!}\\
   & \textbf{erzählen wir!}\\
   & \textbf{erzählt (ihr)!}\\
   & \textbf{erzählen Sie!}\\
   & \vspace{10pt}
  \end{tabular}

    \vspace{10pt}

 \begin{tabular}{rl}
     \multicolumn{2}{l}{\textbf{PERFEKT}}\\
  &  haben + \textbf{erzählt}\\
  \end{tabular}
\hfill
   \begin{tabular}{rl}
   
    \multicolumn{2}{l}{\textbf{FUTURE I} }\\
  &  werden + \textbf{erzählen}\\
 
  \end{tabular}
\hfill
   \begin{tabular}{rl}
      \multicolumn{2}{l}{\textbf{PLUSQUAMPERFEKT}}\\
   & hatten + \textbf{erzählen}\\
  \end{tabular}
  
  
\vspace{10pt}

\begin{tabular}{lll}
\multicolumn{3}{l}{\textbf{MIT PRÄPOSITIONEN}}\\
& \textbf{---} & \\
\end{tabular}
\hfill
\begin{tabular}{lll}
\multicolumn{3}{l}{\textbf{TRENNBARE VERBEN}}\\
& \textbf{---} & \\
\end{tabular}
\vspace{-5pt}
\end{verbentry}

% erziehen (to raise / to educate)
\begin{verbentry}{
  style = {default},
  toc = {erziehen (to raise / to educate)},
  name = {erziehen},
  english = {to raise / to educate},
  type = {irregular},
  auxiliary = {haben}
}
 
    
     \vspace{5pt}

    \hrule
    
    \vspace{5pt}

  \begin{tabular}{rl}
     \multicolumn{2}{l}{\textbf{PRÄSENS} }\\
    ich & \textbf{erziehe}\\
    du & \textbf{erziehst}\\
    er/sie/es & \textbf{erzieht}\\
    wir & \textbf{erziehen}\\
    ihr & \textbf{erzieht}\\
    sie/Sie & \textbf{erziehen}\\
  \end{tabular}
  \hfill
     \begin{tabular}{rl}
    \multicolumn{2}{l}{\textbf{PRÄTERITUM} }\\
    ich & \textbf{erzog}\\
    du & \textbf{erzogst}\\
    er/sie/es & \textbf{erzog}\\
    wir & \textbf{erzogen}\\
    ihr & \textbf{erzogt}\\
    sie/Sie & \textbf{erzogen}\\
  \end{tabular}
  \hfill
  \begin{tabular}{rl}
     \multicolumn{2}{l}{\textbf{IMPERATIV} }\\
   & \textbf{erzieh(e) (du)!}\\
   & \textbf{erziehen wir!}\\
   & \textbf{erzieht (ihr)!}\\
   & \textbf{erziehen Sie!}\\
   & \vspace{10pt}
  \end{tabular}

    \vspace{10pt}

 \begin{tabular}{rl}
     \multicolumn{2}{l}{\textbf{PERFEKT}}\\
   & haben + \textbf{erzogen}\\
  \end{tabular}
\hfill
   \begin{tabular}{rl}
    \multicolumn{2}{l}{\textbf{FUTURE I} }\\
   & werde + \textbf{erziehen}\\
  \end{tabular}
\hfill
   \begin{tabular}{rl}
      \multicolumn{2}{l}{\textbf{PLUSQUAMPERFEKT}}\\
   & hatten + \textbf{erzogen}\\
  \end{tabular}

   \vspace{10pt}

   \begin{tabular}{lll}
      \multicolumn{3}{l}{\textbf{MIT PRÄPOSITIONEN}}\\
& \textbf{erziehen zu} + Dat & (to raise/educate to be)\\
& \textbf{erziehen mit} + Dat & (to raise/educate with)\\
\end{tabular}
  \hfill
   \begin{tabular}{lll}
\multicolumn{3}{l}{\textbf{TRENNBARE VERBEN}}\\
& \textbf{---} & \\
\end{tabular}

  \vspace{-5pt}
\end{verbentry}

% essen (to eat)
\begin{verbentry}{
  style = {acc},
  toc = {essen (to eat)},
  name = {essen},
  english = {to eat},
  type = {irregular, + Akkusativ},
  auxiliary = {haben}
}
 
    
     \vspace{5pt}

    \hrule
    
    \vspace{5pt}

  \begin{tabular}{rl}
     \multicolumn{2}{l}{\textbf{PRÄSENS} }\\
    ich & \textbf{esse}\\
    du & \textbf{isst}\\
    er/sie/es & \textbf{isst}\\
    wir & \textbf{essen}\\
    ihr & \textbf{esst}\\
    sie/Sie & \textbf{essen}\\
  \end{tabular}
  \hfill
     \begin{tabular}{rl}
    \multicolumn{2}{l}{\textbf{PRÄTERITUM} }\\
    ich & \textbf{aß}\\
    du & \textbf{aßest}\\
    er/sie/es & \textbf{aß}\\
    wir & \textbf{aßen}\\
    ihr & \textbf{aßt}\\
    sie/Sie & \textbf{aßen}\\
  \end{tabular}
  \hfill
  \begin{tabular}{rl}
     \multicolumn{2}{l}{\textbf{IMPERATIV} }\\
   & \textbf{iss (du)!}\\
   & \textbf{essen wir!}\\
   & \textbf{esst (ihr)!}\\
   & \textbf{essen Sie!}\\
   & \vspace{10pt}
  \end{tabular}

    \vspace{10pt}

 \begin{tabular}{rl}
     \multicolumn{2}{l}{\textbf{PERFEKT}}\\
  &  haben + \textbf{gegessen}\\
  \end{tabular}
\hfill
   \begin{tabular}{rl}
   
    \multicolumn{2}{l}{\textbf{FUTURE I} }\\
  &  werden + \textbf{essen}\\
 
  \end{tabular}
\hfill
   \begin{tabular}{rl}
      \multicolumn{2}{l}{\textbf{PLUSQUAMPERFEKT}}\\
   & hatten + \textbf{gegessen}\\
  \end{tabular}
  
  
\vspace{10pt}

\begin{tabular}{lll}
\multicolumn{3}{l}{\textbf{MIT PRÄPOSITIONEN}}\\
& \textbf{---} & \\
\end{tabular}
\hfill
\begin{tabular}{lll}
\multicolumn{3}{l}{\textbf{TRENNBARE VERBEN}}\\
& \textbf{---} & \\
\end{tabular}
\vspace{-5pt}
\end{verbentry}
